% p024_a 的电脑重画（动源：电源接两根导轨，中间棒(电阻 R，初速 0)与右边运动后的棒(速度 v，棒上等效电源 BLv 为反电动势)；回路中顺时针电流环与 B；原图左上角“动源”是标题，已在正文中，不重复）
\begin{tikzpicture}[line width=0.7pt, >={Stealth[length=2.2mm]}]
  % 导轨
  \draw (0,2.3) -- (5.0,2.3);
  \draw (0,0) -- (4.6,0);
  % 左：电源（上长板、下短板）
  \draw (0,2.3) -- (0,1.4);
  \draw[line width=1.1pt] (-0.5,1.4) -- (0.5,1.4);
  \draw[line width=1.1pt] (-0.25,1.0) -- (0.25,1.0);
  \draw (0,1.0) -- (0,0);
  % 中间的棒（电阻 R）
  \draw (2.6,-0.7) -- (2.6,3.0);
  \node at (2.6,3.3) {$0$};
  \node[left] at (2.55,1.0) {$R$};
  \draw[->] (2.8,1.7) -- (3.25,1.7);
  % 右边的棒（运动后），棒上有等效电源
  \draw (4.3,-0.7) -- (4.3,3.0);
  \draw[line width=1.1pt] (3.95,1.57) -- (4.65,1.57);
  \draw[line width=1.1pt] (4.12,1.35) -- (4.48,1.35);
  \node at (4.6,3.3) {$v$};
  \node[right] at (4.85,1.3) {$BLv$};
  \node[right] at (4.7,0) {反电动势};
  % 棒与导轨的交点
  \fill (2.6,2.3) circle (1.4pt) (2.6,0) circle (1.4pt) (4.3,2.3) circle (1.4pt) (4.3,0) circle (1.4pt);
  % 左回路中的电流环（顺时针）与 B（向里）
  \draw[postaction={decorate}, decoration={markings, mark=at position 0.727 with {\arrow{>}}}, ->]
    ([shift=(114:0.9)]1.05,1.1) arc (114:-120:0.9);
  \node at (1.2,1.45) {$\times$};
\end{tikzpicture}
